Three components extend the baseline of \cref{sec:proposed}: an $L_1$-adaptive altitude augmentation, a receding-horizon MPC with a penalized tension-ceiling constraint (\cref{sec:mpc-layer}), and a fault-triggered formation-reshape supervisor. The $L_1$ layer compensates parameter deviations $(m_L, k_{\mathrm{eff}})$ from the nominal $q^\star$. The MPC layer enforces a horizon-aware linearized tension ceiling, binding when safe working load dominates. The reshape supervisor reassigns surviving slots to minimise worst-case tension under a symmetric hover model, binding when hover-equilibrium asymmetry rather than dynamic load dominates (P2-D, \cref{rem:reshape-null}). Each preserves H1--H3 (\cref{sec:theorem-preservation}), so \eqref{eq:iss-envelope-explicit} holds with the relevant $\gamma$-gain reduced.

The $L_1$ layer is probed by P2-B (mass-mismatch, FF off), recovering $\sim 40\%$ of the induced altitude sag; MPC by P2-C (tension-ceiling sweep, V4); reshape by P2-D (lemniscate-period sweep, V4). On the demonstrated schedule, peak tension is dominated by the pre-fault dynamic spike; neither MPC nor reshape binds. Each subsection identifies its binding regime; \cref{sec:theorem-preservation} shows enabling a layer outside that regime leaves the envelope unchanged.

\subsection{$L_1$-Adaptive Altitude Augmentation}
\label{sec:l1-augmentation}

The identity $T_i^{\mathrm{ff}} = T_i$ does not cancel matched disturbances entering through parameter deviations $(m_L, k_{\mathrm{eff}})$. The $L_1$ component~\cite{Hovakimyan2010book, CaoHovakimyan2008} estimates this residual online and injects an additive altitude correction $u_{\mathrm{ad}}$ between the outer PD and the QP. In altitude-error coordinates $\tilde z_i \triangleq p_{\mathrm{slot},i,z} - p_{i,z}$, the baseline closed loop is
\begin{equation}
  \ddot{\tilde z}_i \;=\; -K_p^z\,\tilde z_i - K_d^z\,\dot{\tilde z}_i + \delta_m(t),
  \label{eq:l1-channel}
\end{equation}
where $\delta_m(t)$ is the residual $r_i(t)$ of \cref{lem:tension-cancel} augmented by parameter-mismatch components, with $|\delta_m(t)| \le \delta_{m,\max}$.

The standard $L_1$ architecture~\cite{Hovakimyan2010book}---Hurwitz reference model, state predictor, projected gradient on $\hat\delta_m$, low-pass filter---is adopted verbatim with $A_m$, $P_v$ from \cref{lem:altitude-iss}. The adaptive update
\begin{align}
\label{eq:l1-adaptive}
  \dot{\hat\delta}_m
    &\;=\; -\Gamma\,\operatorname{Proj}\!\bigl(
      \hat\delta_m,\;\vect b^\top P_v\,\vect{\tilde x}\bigr),\\
  \hat\delta_m &\in [\delta_{m,\min}, \delta_{m,\max}],\nonumber
\end{align}
runs at $T_s = \SI{2e-4}{\second}$; the filtered $u_{\mathrm{ad}}$ is added to the altitude component of $\vect a_{\mathrm{target}}$ \eqref{eq:a-target}, at the injection point labelled in \cref{fig:controller-cascade}. The cutoff $\omega_c = \SI{25}{\radian\per\second}$ separates the estimator from the reference natural frequency $\omega_n^z = \sqrt{K_p^z} = \SI{10}{\radian\per\second}$ by a factor of $2.5$. Contractivity of the scalar Euler-discretized update gives the following bound.

\begin{proposition}[Scalar $L_1$ admissible-window bound]
\label{prop:gamma-star}
For the Euler-discretized adaptive-law update \eqref{eq:l1-adaptive} running at step $T_s$, the scalar proxy of the discrete adaptation map yields two necessary conditions:
\begin{enumerate}[label=(\roman*),leftmargin=2em]
  \item \emph{(Discretization stability)}
    \begin{equation}
      \Gamma \;<\; \Gamma^\star
        \;\triangleq\; \frac{2}{T_s\,p_{22}};
      \label{eq:gamma-star}
    \end{equation}
  \item \emph{(Bandwidth coverage)} the continuous-time adaptation rate
    must lie within the filter passband,
    \begin{equation}
      \Gamma \;>\; \Gamma_{\min}
        \;\triangleq\; \frac{\omega_c}{p_{22}}.
      \label{eq:gamma-min}
    \end{equation}
\end{enumerate}
For $T_s = \SI{2e-4}{\second}$, $p_{22} = 0.0210$, $\omega_c = \SI{25}{\radian\per\second}$:
\begin{equation}
  \Gamma_{\min} \;\approx\; 1{,}190
  \;<\;
  \underbrace{\Gamma = 2{,}000}_{\text{operating point}}
  \;<\;
  \Gamma^\star \;\approx\; 4.75\times 10^5.
  \label{eq:gamma-window}
\end{equation}
\end{proposition}

\begin{proof}
Isolating the scalar $\hat\delta_m$-update in the absence of predictor error, projection, and filter dynamics gives the one-step map $\hat\delta_m \mapsto (1-T_s\Gamma p_{22})\hat\delta_m + \mathrm{const}$, with scalar pole $z^\star = 1 - T_s\Gamma p_{22}$.

\smallskip
\noindent\emph{(i) Stability.} $|z^\star| < 1$ requires $T_s\Gamma p_{22} \in (0,2)$, which gives \eqref{eq:gamma-star}.

\smallskip
\noindent\emph{(ii) Bandwidth.} The equivalent continuous-time adaptation rate is $\gamma_{\mathrm{ad}} = -\ln(z^\star)/T_s$. Since $T_s\Gamma p_{22} = 8.4\times10^{-3} \ll 1$ at the operating point, $\gamma_{\mathrm{ad}} \approx \Gamma p_{22}$ with error $< 0.4\%$. The low-pass filter with bandwidth $\omega_c$ attenuates disturbances above $\omega_c$ before the adaptation can correct them; avoiding this requires $\gamma_{\mathrm{ad}} \ge \omega_c$, i.e.\ $\Gamma p_{22} \ge \omega_c$, which gives \eqref{eq:gamma-min}.
\end{proof}

\begin{remark}[Design interpretation of \eqref{eq:gamma-window}]
Both bounds operate on the scalar proxy; the full discrete interconnection (predictor error, projection, filter dynamics) is handled by composite Lyapunov analysis following~\cite{Hovakimyan2010book}. Within the proxy, \eqref{eq:gamma-window} partitions the gain axis into three regimes: below $\Gamma_{\min}$ the adaptation is too slow to track disturbances within the filter passband; above $\Gamma^\star$ the Euler integration becomes contractive only in the unstable direction. The canonical $\Gamma = 2{,}000$ is $1.68\times\Gamma_{\min}$ — chosen on the performance side of the window where the adaptation rate is $68\%$ faster than the filter bandwidth. The $238\times$ gap to $\Gamma^\star$ is a deliberately maintained discretization safety margin, not a sign of conservatism; operating near $\Gamma_{\min}$ (not near $\Gamma^\star$) keeps the adaptation rate well separated from the sampling frequency and avoids exciting unmodelled predictor-error cross-coupling. Composite Lyapunov analysis yields the standard $L_1$ ISS guarantee w.r.t.\ the unmatched residual $d_\perp$; matched components are rejected up to $\omega_c$. Empirical validation (P2-B) is in \cref{sec:VI-G}.
\end{remark}

\subsection{Receding-Horizon MPC with a Penalized Tension-Ceiling Constraint}
\label{sec:mpc-layer}

\paragraph{Role in the architecture.} The single-step QP \eqref{eq:qp} projects the current-tick target onto the actuator envelope and cannot foresee a tension spike one or more ticks ahead. When safe working load is the dominant requirement, we replace it with an $N_p$-step receding-horizon MPC carrying a penalized linearized tension-ceiling constraint at every horizon step. The constraint is soft: a non-negative slack $\vect s \ge \vect 0$ is added to the ceiling inequality and penalized with weight $w_s = 10^4$, trading unconditional feasibility against ceiling excursions when the envelope is tight. The MPC reduces to the baseline QP when the ceiling is loose; it becomes binding when the predicted tension approaches the ceiling.

\paragraph{Prediction model and tick QP.} Let $\vect x[k] \triangleq (\vect p_{\mathrm{slot},i} - \vect p_i;\,\vect v_{\mathrm{slot},i} - \vect v_i) \in \R^6$ be the error state at $t_k = k\,\Delta t_{\mathrm{MPC}}$ and $\vect u[k] \in \R^3$ the commanded acceleration. ZOH discretization gives $\vect x[k{+}1] = A\,\vect x[k] + B\,\vect u[k]$ \eqref{eq:mpc-AB}, and horizon stacking yields $\mathcal{X} = \Phi\,\vect x[0] + \Omega\,\mathcal{U}$. Linearising $T_i = k_{\mathrm{eff}}\max(0,\|\vect p_i - \vect p_L\| - L_{\mathrm{eff}})$ around the taut chord $\vect c[k]$ with unit direction $\hat{\vect n}[k]$ gives the horizon-$j$ prediction
\begin{align}
\label{eq:T-linearized}
  T_i[k+j] \;\approx\;
    T_i[k]
    &+ j\,\Delta t_{\mathrm{MPC}}\,\dot T_i^{\mathrm{filt}}[k]\\
    &+ \vect c_j^\top (A^j - I)\,\vect x[k]+ \vect c_j^\top \Omega_{\mathrm{row}\,j}\,\mathcal{U},\nonumber
\end{align}
with $\vect c_j = (-k_{\mathrm{eff}}\hat{\vect n}[k];\,\vect 0)$. The linearization error is bounded by $k_{\mathrm{eff}} E^2/(2 d_{\mathrm{nom}})$ in position-error magnitude $E$ and nominal chord length $d_{\mathrm{nom}}$. The state-space recursion is
\begin{align}
\label{eq:mpc-AB}
  \vect x[k+1] = A\vect x[k] + B\vect u[k],&\quad A = \begin{bmatrix} I_3 & \Delta t_{\mathrm{MPC}} I_3 \\ \mat 0 & I_3\end{bmatrix},\\
  B = \begin{bmatrix}-\tfrac{1}{2}\Delta t_{\mathrm{MPC}}^2 I_3 \\ -\Delta t_{\mathrm{MPC}} I_3\end{bmatrix}.&\nonumber
\end{align}

The tick QP combines an outer-loop tracking objective with the slack-penalised tension constraint:
\begin{align}
\label{eq:mpc-qp}
  \min_{\mathcal{U},\,\vect s \ge \vect 0}
    \tfrac{1}{2}\mathcal{U}^\top H \mathcal{U}
    + \vect f^\top \mathcal{U}
    + w_s\,\vect 1^\top \vect s\\
  \text{s.t.}
  \quad
  G\,\mathcal{U} \le \vect d_T(\vect x[k]) + \vect s,\nonumber
\end{align}
where $G$ stacks $\vect c_j^\top \Omega_{\mathrm{row}\,j}$ and $\vect d_T(\vect x[k])$ collects the state-dependent right-hand sides. The QP is convex and solved by a single warm-started OSQP call per tick~\cite{stellato2020osqp}; only $\vect u[0]$ is applied. A discrete-Riccati terminal cost yields LQR gain with $\rho_{\mathrm{spec}}(A - BK) \approx 0.969$.

\begin{corollary}[Unconditional feasibility]
\label{cor:mpc-recfeas}
\eqref{eq:mpc-qp} is feasible at every tick: $\mathcal{U} = \vect 0$ with $\vect s[k{+}j] = \max(0, -[\vect d_T(\vect x[k])]_{j+1})$ is always admissible.
\end{corollary}

\begin{proposition}[Slack-activation regime]
\label{prop:mpc-slack}
Let $\|q - q^\star\|_\infty \le \bar\varepsilon_q$ with $q = (m_L, k_{\mathrm{eff}})$ and $\vect d_T$ $L_q$-Lipschitz in $q$. Under D1 and an initial tick where the unconstrained QP satisfies the ceiling, $\|\vect s[k]\|_\infty \le L_q \bar\varepsilon_q + O(\bar\varepsilon_q^2)$, i.e.\ slack activation scales as $O(\bar\varepsilon_q)$.
\end{proposition}

\begin{proof}[Sketch]
$\vect d_T$ depends linearly on $q$ through $k_{\mathrm{eff}}$ in $\vect c_j$; Taylor expansion around $q^\star$ gives $\vect d_T(q) = \vect d_T(q^\star) + D_q \vect d_T (q-q^\star) + O(\|\cdot\|^2)$. Since the ceiling holds without slack at $q^\star$, the minimum slack restoring feasibility at $q \ne q^\star$ is $\le L_q \|q-q^\star\|_\infty + O(\|\cdot\|^2)$. D1 preserves linearization on each inter-fault window.
\end{proof}

Body-to-body cable geometry requires effective rest length $L_{\mathrm{eff}}^{\mathrm{body}} = \sqrt{r_f^2 + \rho_{\mathrm{geo}}^2} - \delta^\star \approx \SI{1.448}{\metre}$, set so the linearized elastic force at hover matches the simulator's measured \SI{8}{\newton}. The empirical activation regime (P2-C) is in \cref{sec:VI-H}; the layer is non-binding on the demonstrated family because peak tension is set by the pre-fault dynamic spike rather than the quasi-static load the ceiling targets. MPC parameters: $N_p = 5$ (10 at evaluation), $\Delta t_{\mathrm{MPC}} = \SI{0.01}{\second}$, $q_p = 100$, $q_v = 10$, $r = 0.02$, $w_s = 10^4$, $T_{\max} = \SI{100}{\newton}$, $k_{\mathrm{eff}} = \SI{2778}{\newton\per\metre}$.

\subsection{Fault-Triggered Formation-Reshape Supervisor}
\label{sec:reshape-supervisor}

After a severance, the survivors retain their original angular slots, leaving the hover-equilibrium tension distribution asymmetric. This asymmetry is the binding tension load only when the trajectory is sufficiently low-bandwidth that the dynamic centripetal component does not dominate (P2-D, \cref{rem:reshape-null}). The reshape supervisor reassigns survivors to an equiangular configuration through a $C^2$ quintic smoothstep over $T_{\mathrm{trans}} = \SI{5}{\second}$; once the new layout is reached, the supervisor latches and the architecture returns to pre-fault flow. Unlike the baseline, this is an \emph{explicitly supervised} extension: each drone runs a local tension-latch detector ($T_{\mathrm{fault}} = \SI{0.5}{\newton}$, $\tau_{\mathrm{det}} = \SI{100}{\milli\second}$), and a per-fault broadcast of at most $\lceil \log_2 N\rceil$ bits distributes the fault index. The no-detection, no-communication guarantees of \cref{thm:hybrid-stability} do not extend here.

\begin{theorem}[Globally tension-optimal reshape, $N = 4$]
\label{thm:reshape}
Let $N = 4$ with nominal slots $\phi_i = \pi i / 2$ and let cable $i^\star$ sever. Under the symmetric-load-sharing model ($T_k = (m_L g)/(M\cos\beta)(1 + \epsilon_k)$ with angular residual $\epsilon_k$), the worst-case-tension-minimising assignment over the survivors is: the drone diametrically opposite $i^\star$ stays at its nominal angle, and the two adjacent drones rotate $\pi/6$ toward the gap. The resulting equiangular $120^\circ$ spacing minimises $\max_i T_i$ over all continuous assignments respecting hover equilibrium, with a closed-form worst-case reduction of $25.8\%$.
\end{theorem}

\begin{proof}[Sketch]
By reflection symmetry about the gap axis, any optimum has a symmetric copy; with the opposite drone at $\Delta\phi = 0$ and a common offset for the adjacent pair, the residuals are piecewise quadratic in $\Delta\phi$ and the minimax occurs at $\Delta\phi = \pi/6$. The endpoints $\Delta\phi \in \{0, \pi/4\}$ are dominated.
\end{proof}

The $V_4 \to V_6$ dual-fault scenarios traverse $N = 5 \to 4 \to 3$: \cref{thm:reshape} applies directly at the second fault; at the first, the surviving four are assigned to $90^\circ$ spacing by the same minimax argument. A formal arbitrary-$N$ generalisation is follow-up work. The empirical regime is reported in \cref{sec:VI-H}: NF3 returns $0.0\%$ benefit across $T_{\mathrm{ref}} \in \{6, 8, 10, 12\}$~s because the dynamic centripetal load dominates the quasi-static-hover target the assignment optimises against.

\subsection{Baseline Theorem Preservation and Admissibility
Regime}
\label{sec:theorem-preservation}

Each component modifies the closed loop through a disjoint channel and preserves H1--H3 of \cref{thm:hybrid-stability} on $\Omega = \Omega_\tau^{\mathrm{dwell}} \cap \Omega_{\mathrm{QP, active}}$. The $L_1$ layer adds $u_{\mathrm{ad}}$ to the altitude component of $\vect a_{\mathrm{target}}$ before the QP, leaving the QP active set and the identity unchanged; H1, H3 hold verbatim, H2 is controller-independent, and the layer reduces $\gamma_q$ in \eqref{eq:iss-envelope-explicit} by the composite $V_{L_1}$ contraction. The MPC layer replaces \eqref{eq:qp} with \eqref{eq:mpc-qp}; \cref{cor:mpc-recfeas} guarantees unconditional feasibility, \cref{prop:mpc-slack} bounds the slack by the parameter mismatch, and the closed loop reduces to the baseline in the unconstrained regime. The reshape supervisor injects a $C^2$-bounded slot perturbation through the override hook; H1, H3 hold provided the smoothstep keeps the perturbation in the admissible-reference set, and H2 reduces to a latching rule firing at most once per dwell period. \eqref{eq:iss-envelope-explicit} therefore remains valid under any combination of enabled components; cross-layer interactions are evaluated empirically.

\paragraph{Admissibility regime and out-of-scope conditions.} The admissibility domain, $\rho$-gap analysis, and conditional certificate are as in \cref{rem:rho-gap,sec:stability-setup}; the three extension layers do not alter these bounds.

Three regimes lie outside the analysis domain. \emph{Rapid multi-fault} ($\tau_d < \tau_{\mathrm{pend}}$) violates H2 and voids the contraction of \cref{lem:dwell-decay}. \emph{Long slack excursions} (slack-run $> \SI{40}{\milli\second}$ or duty $> 2.5\%$) violate H1 and exit $\Omega_\tau^{\mathrm{dwell}}$. \emph{QP fallback}, in which infeasibility forces the component-wise clamp, violates H3 and breaks D1.

A second class lies outside the architectural category: faults that do not arrive on the rope-tension channel. \emph{Drone-side faults} (rotor loss, IMU failure, brownout) do not self-announce through $T_i$; \emph{soft cable failures} (fraying, attachment slip) produce a graded rather than discrete tension signature; \emph{tension-sensor failures} disable the announcement itself; \emph{payload-attitude dynamics} require rigid-body coupling beyond the point-mass model. These are not gaps in the campaign; they are out-of-category for the architectural thesis.

Future work inside the category includes a multi-fault extension of the dwell argument under controlled fault-schedule density (closing the rapid-multi-fault and QP-fallback regimes analytically); a reshape extension whose objective accounts for the dynamic centripetal load, since the aggressive-period probe (E8, $T_{\mathrm{ref}}=\SI{6}{\second}$) shows the quasi-static-hover prediction of \cref{thm:reshape} does not bind on the present reference family; and tube-MPC tightening that folds $\Gamma_f e^{-\lambda(t - t_k^\star)}$ into the constraint right-hand side for spike-dominated regimes.
